\documentclass[10pt]{article}
\usepackage[T1]{fontenc}
\usepackage[utf8]{inputenc}
\usepackage[french]{babel}
\usepackage{amsmath,amssymb,mathrsfs,mathpazo}
\usepackage{geometry}
\geometry{a4paper,margin=2cm}
\usepackage{graphicx,array,tabularx,tabularray,colortbl,multirow,multicol,booktabs,enumitem,pifont,ulem,fancybox,diagbox,cancel}
\usepackage{pstricks,pstricks-add,pst-plot,pst-tree,pst-node,pst-text,pst-func,pst-math,pst-bspline,pst-3dplot,pst-eucl,pst-solides3d}
\let\pscalculate\undefined
\usepackage{pst-bezier}
\usepackage{tikz,pgfplots,tkz-tab}
\usetikzlibrary{arrows,arrows.meta,patterns,calc,positioning,trees,decorations.markings,shapes.geometric,shapes.misc,shapes.arrows,intersections,decorations.pathmorphing,decorations.pathreplacing,backgrounds,shadows,babel}
\pgfplotsset{compat=1.18}
\usepgfplotslibrary{fillbetween}
\usepackage[np]{numprint}
\usepackage{listings,makecell,pas-tableur,twoopt}
\let\barre\undefined
\usepackage[pstricks]{tablvar}
\definecolor{darkgreen}{RGB}{0,100,0}
\providecommand{\red}{\color{red}}
\providecommand{\blue}{\color{blue}}
\providecommand{\green}{\color{green}}
\providecommand{\black}{\color{black}}
\providecommand{\white}{\color{white}}
\providecommand{\gray}{\color{gray}}
\providecommand{\cyan}{\color{cyan}}
\providecommand{\magenta}{\color{magenta}}
\providecommand{\yellow}{\color{yellow}}
\providecommand{\violet}{\color{violet}}
\providecommand{\orange}{\color{orange}}
\providecommand{\euro}{\texteuro}
\providecommand{\EUR}{\texteuro}
\providecommand{\ieme}{\textsuperscript{e}}
\providecommand{\degres}{\ensuremath{^\circ}}
\providecommand{\degree}{\ensuremath{^\circ}}
\providecommand{\vv}[1]{\overrightarrow{#1}}
\providecommand{\Coord}[1]{#1}
\providecommand{\decosix}{}
\providecommand{\starredbullet}{\ensuremath{\bullet}}
\pagestyle{empty}
\setlength{\parindent}{0pt}
\usepackage{adjustbox}\DeclareUnicodeCharacter{2236}{:}\DeclareUnicodeCharacter{2190}{\ensuremath{\leftarrow}}
\providecommand{\R}{\mathbb{R}}
\providecommand{\N}{\mathbb{N}}
\providecommand{\Z}{\mathbb{Z}}
\providecommand{\Q}{\mathbb{Q}}
\providecommand{\D}{\mathbb{D}}
\providecommand{\C}{\mathbb{C}}
\providecommand{\e}{{\mathrm{e}}}
\providecommand{\vect}[1]{\overrightarrow{#1}}
\providecommand{\vectt}[1]{\overrightarrow{\mathrm{#1}}}
\providecommand{\vvt}[1]{\overrightarrow{\mathrm{#1}}}
\providecommand{\np}[2][]{{#2~#1}}
\providecommand{\vv}[1]{{\overrightarrow{#1}}}
\renewcommand{\d}{{\mathrm{d}}}
\providecommand{\degres}{{^\circ}}
\providecommand{\degre}{{^\circ}}
\providecommand{\blue}{\color{blue}}
\providecommand{\red}{\color{red}}
\providecommand{\enskip}{\quad}
\providecommand{\llbracket}{\lbrack\!\lbrack}
\providecommand{\rrbracket}{\rbrack\!\rbrack}
\providecommand{\euro}{{\text{€}}}
\providecommand{\no}{{\text{n°}}}
\providecommand{\interval}[2]{{[#1;#2]}}
\providecommand{\intervFO}[2]{{[#1;#2[}}
\providecommand{\dotfill}{\ldots}

\providecommand{\ds}{\displaystyle}
\providecommand{\textding}[1]{#1}
\providecommand{\Oij}{\(\mathrm{O};\vec{\imath},\vec{\jmath}\)}
\providecommand{\Oijk}{\(\mathrm{O};\vec{\imath},\vec{\jmath},\vec{k}\)}
\providecommand{\Ouv}{\(\mathrm{O};\vec{u},\vec{v}\)}
\providecommand{\intervFF}[2]{\ensuremath{[#1;#2]}}
\renewcommand{\euro}{\ifmmode\text{\texteuro}\else\texteuro\fi}
\DeclareUnicodeCharacter{00B0}{\ifmmode^{\circ}\else\textdegree\fi}
\DeclareUnicodeCharacter{20AC}{\ifmmode\text{\texteuro}\else\texteuro\fi}
% Énoncés : Éducation nationale. Transcription et corrigés : APMEP, Denis Vergès et contributeurs. Découpage et classement : CoopMaths, Sébastien Lozano, Éric Elter et contributeurs. Adaptation de la présentation : Maths Explorer.
% CC BY-SA : https://coopmaths.fr/www/about/
% https://coopmaths.fr/alea/static/bac/2024/tex/bac_2024_11_sujet1_amsud_1.tex
\begin{document}
\noindent\textit{APMEP / CoopMaths. CC BY-SA. Présentation adaptée par Maths Explorer.}\par\medskip
\medskip

\textbf{PARTIE A}

\medskip

On considère l'équation différentielle 
\[(E) : \quad y' +\dfrac14y = 20\e^{-\frac14 x},\]
d'inconnue $y$, fonction définie et dérivable sur l'intervalle $[0~;~ +\infty[$.

\medskip

\begin{enumerate}
\item Déterminer la valeur du réel $a$ tel que la fonction $g$ définie sur l'intervalle $[0~;~ +\infty[$ par $g(x) = ax\e^{-\frac14 x}$ soit une solution particulière de l'équation différentielle $(E)$.

\item On considère l'équation différentielle 
\[(E') : \quad y' +\dfrac14 y = 0,\]
d'inconnue $y$, fonction définie et dérivable sur l'intervalle $[0~;~ +\infty[$.

Déterminer les solutions de l'équation différentielle $(E')$.
\item En déduire les solutions de l'équation différentielle $(E)$.
\item Déterminer la solution $f$ de l'équation différentielle $(E)$ telle que $f(0) = 8$.
\end{enumerate}

\medskip

\textbf{PARTIE B}

\medskip

On considère la fonction $f$ définie sur l'intervalle $[0~;~ +\infty[$ par 
\[f(x) = (20x + 8)\e^{-\frac14 x}.\]

On admet que la fonction $f$ est dérivable sur l'intervalle $[0~;~ +\infty[$ et on note $f'$ sa fonction dérivée sur l'intervalle $[0~;~ +\infty[$. De plus, on admet que $\displaystyle\lim_{x \to + \infty}f(x) = 0$.

\medskip

\begin{enumerate}
\item 
	\begin{enumerate}
		\item Justifier que, pour tout réel $x$ positif, 
\[f'(x) = (18 - 5x)\e^{-\frac14 x}.\]

		\item En déduire le tableau de variations de la fonction $f$. On précisera la valeur exacte du maximum de la fonction $f$ sur l'intervalle $[0~;~ +\infty[$.
	\end{enumerate}
\item Dans cette question on s'intéresse à l'équation $f(x) = 8$.
	\begin{enumerate}
		\item Justifier que l'équation $f(x) = 8$ admet une unique solution, notée $\alpha$, dans l'intervalle [14~;~15].
		\item Recopier et compléter le tableau ci-dessous en faisant tourner étape par étape la fonction solution\_equation ci-contre, écrite en langage Python

\begin{minipage}{0.61\linewidth}
\begin{tabularx}{\linewidth}{|m{2.25cm}|c|*{4}{>{\centering \arraybackslash}X|}}\hline
	$a$		&14&&&&\\ \hline
	$b$		&15&&&&\\ \hline
	$b - a$	&1&&&&\\ \hline
	$m$		&14,5&&&&\cellcolor[gray]{0.2}\\ \hline
Condition $f(m) > 8$	&FAUX	&	&	&	&\cellcolor[gray]{0.2}\\ \hline
\end{tabularx}
\end{minipage}\hfill
\begin{minipage}{0.36\linewidth}
{\footnotesize\begin{tabular}{|l|}\hline
from math import exp \\
def f(x):\\
\qquad return (20*x+8)*exp(-1/4*x)\\
~\\
def solution\_equation():\\
\qquad a,b = 14,15\\
\qquad while b-a > 0.1 :\\
\qquad\begin{tabular}{|l}
\quad m = (a+b)/2\\
\quad if f(m) > 8:\\
\quad|\quad a = m\\
\quad else:\\
\quad |\quad  b = m\\
\end{tabular}\\
\qquad return a,b\\ \hline
\end{tabular}}

\end{minipage}
		\item Quel est l'objectif de la fonction solution\_equation
dans le contexte de la question ?
	\end{enumerate}
\end{enumerate}

\bigskip

\end{document}
